\chapter{Segments, Indexing, Retention, and Truncation}
As a single log file grows, restart scans become longer; deleting old data from the middle would otherwise require rewriting everything after it. A partition log is therefore made of consecutive segment files: each segment contains a range of offsets, the active segment receives appends, and closed segments are no longer appended to. A byte position within a segment counts file bytes, not partition offsets. A sparse index stores only some offset-to-byte-position hints and can be rebuilt from the main log; it speeds lookup but is not record truth. Retention deletes a complete old segment prefix, while truncation deletes a tail after a nextOffset—their directions and purposes differ.

\begin{figure}[ht]
\centering
\begin{tikzpicture}[x=0.8mm,y=1mm,font=\small]
\draw[thick,rounded corners] (0,0) rectangle (57,22); \node at (28.5,16) {segment base 000};
\node at (28.5,8) {offset 0--3};
\draw[thick,rounded corners] (65,0) rectangle (122,22); \node at (93.5,16) {segment base 004};
\node at (93.5,8) {offset 4--7};
\draw[thick,rounded corners,fill=blue!5] (130,0) rectangle (187,22); \node at (158.5,16) {active base 008};
\node at (158.5,8) {LEO=8};
\draw[-{Stealth},thick] (28, -4) -- (28, -13); \node[below] at (28,-13) {index floor(2)=(0,pos)};
\draw[-{Stealth},thick] (94, -4) -- (94, -13); \node[below] at (94,-13) {index floor(6)=(4,pos)};
\end{tikzpicture}
\caption{Logical ranges in a segmented log and the sparse index within each segment. The base offset in a segment filename is its logical starting point.}
\end{figure}

\lessonsection{5}{Sparse Index}
\goals{Build a persistent, rebuildable sparse mapping from offsets to byte positions.}{Understand strictly increasing observations, floor lookup, and scanning the primary log.}
\background{An index entry is two big-endian int64 values, \pathcode{offset, position}, for 16 bytes total. Rather than index every record, persist a point every intervalRecords records according to the logical distance from the segment base; always persist the first record. floor(target) chooses the greatest indexed offset not larger than the target, so a later scan can begin at a known record boundary. Even with no index entry, scanning safely from \pathcode{(baseOffset,0)} is possible. A missing or damaged index can be rebuilt from the primary log, but an index cannot repair a damaged log.}
\providededit{The \method{SparseIndex} path, baseOffset, intervalRecords, persisted file, loading, and close lifecycle are provided; \method{SegmentLog} can rescan the primary log.}{Implement \method{public synchronized void add(long offset,long position)}, \method{public synchronized IndexEntry floor(long offset)}, and \method{public synchronized void rebuild(Path logFile,long requestedBaseOffset)} in \pathcode{exercises/src/main/java/io/simplekafka/storage/SparseIndex.java}.}
\subsubsection*{Prerequisites}
Complete \stepref{4}{crash recovery} and know how to validate complete records, CRCs, and continuous offsets in the main log. An index entry's position is a byte count within the segment file; its offset is a record number in the partition. They cannot be interchanged.
\subsubsection*{Inputs, outputs, and state}
\begin{itemize}
\item \method{add} receives a record's observed logical offset and its byte position within the segment. The first observation must be (baseOffset,0); every later observation must strictly increase both offset and position, and position must be nonnegative. An observation that is not persisted still participates in the next ordering check.
\item Persist an observation only when \pathcode{(offset-baseOffset) mod intervalRecords = 0}; the first (baseOffset,0) is always persisted. A rejected observation reports \method{INVALID_REQUEST} and leaves the index file and ordering state unchanged. A closed index or I/O failure reports \method{STORAGE_ERROR}; \method{floor} on a closed index also reports \method{STORAGE_ERROR}.
\item \method{floor} returns the greatest indexed entry at or before target; if none qualifies (including when target is below base), return (baseOffset,0). \method{rebuild} receives the main-log path and required base, clears old entries, scans and validates the log, and recreates all points that should be persisted; it does not change log contents. An invalid path or mismatched base reports \method{INVALID_REQUEST}; corrupt log contents report \method{CORRUPT_RECORD}; a closed index or underlying I/O failure reports \method{STORAGE_ERROR}.
\end{itemize}
\lessoncontract{Each index-file entry stores an 8-byte offset followed by an 8-byte position, both big-endian. Use every observation to validate strict increase; calculate persistence frequency from the offset's distance from base, not from the number of add calls or a fixed byte interval. Missing, truncated, stale, or invalid indexes must be rebuildable from the current log; rebuilding must not make an invalid log appear valid.}
\subsubsection*{Hand calculation: index points and floor}
Assume baseOffset=8 and intervalRecords=2. The segment contains offsets 8, 9, 10, 11; each record has a null key and a 1-byte value, so each occupies 33 bytes and starts at positions 0, 33, 66, 99. Observe all four in order: the first offset 8 writes (8,0); offset 9 is checked for order but not persisted; offset 10 is two records from base, so it writes (10,66); offset 11 is checked but not persisted. The index file contains exactly two 16-byte entries; floor(11)=(10,66), floor(9)=(8,0), and the no-candidate convention makes floor(7)=(8,0) too. If the index is deleted, scanning the four intact log records reconstructs the same two entries.
\expected{\method{Step05Test}: \pathcode{writesBigEndianEntriesAtIntervalAndFloorsSafely} checks 16-byte entries, interval, and floor; \pathcode{supportsNonzeroBasesAndReopensPersistedPoints} checks nonzero bases and reopening; \pathcode{rejectedObservationsLeaveTheIndexAndOrderingStateUnchanged} checks state preservation; \pathcode{rebuildsMissingTruncatedAndStaleIndexesFromTheCurrentLog} and \pathcode{rebuildRejectsMalformedLogsWithoutChangingTheirBytes} check index rebuilding and rejection of malformed logs.}
\pitfalls{floor is not “any nearest entry”: it must not be after the target. An index position must point to a record start or a reader may interpret payload bytes as a header. The index can be discarded and rebuilt; the main log cannot be replaced with a guess.}
\lessonhints{Write the candidate condition offset\(\le\)target, then define the result when there are no candidates: (base,0).}{Separate the ordering state for every observation from whether that point is persisted; the interval is based on offset-base.}{Delete the index file, create SparseIndex again, and observe rebuilding through a real log read rather than checking only an in-memory floor.}
\lessoncommand{5}
\lessonquestions{Why must the primary log still be scanned after finding a floor entry?}{Why can a bad index be rebuilt while a bad primary log cannot use the same recovery strategy?}
\paragraph*{Self-check explanation}
The index is only a search start that does not pass the target; it usually has no exact point for the target, so the records still need to be validated and scanned. An index is redundant guidance that can be recalculated from the record facts; corrupted log contents have no other trusted source from which to infer them.

\lessonsection{6}{Segmented Log}
\goals{Combine single-segment logs into a capacity-rolled partition log and make reads continuous across segments.}{Understand active segments, complete-record capacity, index hints, and global LEO.}
\background{\method{PartitionLog} holds multiple \method{SegmentLog} instances in baseOffset order. An active segment is the final segment that can still receive appends; if it cannot fit the next whole record, close it and create a new segment at the partition's current LEO. Capacity is based on complete encoded record bytes, not value length. When the segment is empty, one legal record larger than segmentBytes may occupy it; otherwise that record could never be appended. A cross-segment read finds the segment containing the target offset, uses that segment's index floor as a hint, then scans forward; maxRecords/maxBytes are the total budgets for the whole read and do not reset at a segment boundary.}
\providededit{Each segment's create/open/recover, sparse index, file discovery, and close path are provided. The \method{PartitionLog} skeleton supplies \method{mutationVersion()} and private \method{markMutation()}.}{Implement \method{public synchronized AppendResult append(List<RecordData> records)} and \method{public List<LogRecord> read(long offset,int maxRecords,int maxBytes)} in \pathcode{exercises/src/main/java/io/simplekafka/storage/PartitionLog.java}. The \method{PartitionLog(Path,long,int)} constructor is provided.}
\subsubsection*{Prerequisites}
Complete \stepref{5}{sparse index}; understand SegmentLog append/read and the index's floor hint. Segment discovery recognizes only regular files named with 20 digits and the \pathcode{.log} suffix; the paired \pathcode{.index} is index data. Other files and directories are not segments to clean up.
\subsubsection*{Inputs, outputs, and state}
\begin{itemize}
\item append takes an ordered batch and returns one global half-open range \pathcode{[firstOffset,nextOffset)}; offsets remain continuous across segment boundaries. An empty/invalid batch or offset overflow reports \method{INVALID_REQUEST} and must not create a segment first. A closed log or underlying I/O failure reports \method{STORAGE_ERROR}; failure after a valid append starts does not promise to roll back files.
\item read accepts an offset in the retained range \pathcode{[logStartOffset,LEO]}; LEO returns empty. maxRecords/maxBytes must be positive. Results are ordered and continuous, and share total record-count and complete encoded-byte budgets across all segments. Out-of-range offset is \method{OFFSET_OUT_OF_RANGE}; a segment gap or corrupt record is \method{CORRUPT_RECORD}; I/O/closed-log failure is \method{STORAGE_ERROR}. If the first record exceeds the byte budget, return empty.
\item Each fully prevalidated append batch increments mutationVersion on this open object once before its first file write or roll; a batch spanning segments still increments once. Invalid input does not increment it. This in-memory mutation generation is not persisted; an I/O failure after writing starts also means an earlier recovery proof is no longer trustworthy. LEO and segment files change on successful append; reads do not change them.
\end{itemize}
\lessoncontract{Roll only when a nonempty active segment cannot fit the next complete record; an empty segment may hold one legal oversized record on its own. Precheck the full batch before first create/roll so invalid input cannot leave an empty segment. Reopening scans registered segments and recovers LEO; unfamiliar files and directories must remain untouched.}
\subsubsection*{Hand calculation: segment rolling and a cross-segment budget}
Case A: segmentBytes=64, with three ordinary records each having a null key and empty value, for 32 bytes each. Offsets 0 and 1 exactly fill the base-0 segment at 64 bytes; before offset 2, roll to base 2, whose size is 32. One append returns [0,3), bases=[0,2], sizes=[64,32]. Case B: three records in order occupy 32, 132, and 32 bytes, with segmentBytes still 64. The 132-byte record cannot fit in the existing nonempty segment, so it occupies a segment by itself at base 1; the following 32-byte record starts a new segment at base 2. Segment sizes are [32,132,32]. A segment may exceed its target only because one legal record itself is larger and the segment was empty.
\par For case A, read(0,3,64) returns only 0 and 1; a 96-byte budget returns 0, 1, and 2. When moving into the second segment, the remaining budget is 32; it must not reset to 64.
\expected{\method{Step06Test}: \pathcode{rollsAtExactByteBoundariesAndWritesIndependentRecords} checks segment capacity; \pathcode{appliesReadBudgetsAcrossSegmentsAndValidatesBoundsAndLimits} checks cross-segment budgets; \pathcode{invalidBatchDoesNotMutateAnExistingMultiSegmentLog} checks no directory mutation for invalid batches; \pathcode{acceptsAnOversizedRecordBetweenRecordsInOneBatchAndReopensEverySegment} checks a legal oversized record alone; \pathcode{recoversOnlyAnIncompleteActiveTailAndRebuildsIndexesOnReopen} and \pathcode{refusesCompleteCrcCorruptionAndSegmentBaseGapsWithoutChangingLogs} check recovery boundaries; \pathcode{concurrentBatchesStayContiguousAndSurviveReopen} checks continuity after concurrent appends and reopen.}
\pitfalls{Counting only payload misses the length, CRC, and fixed fields; resetting maxBytes for each segment exceeds the caller's total budget. A single-segment append does not roll segments—rolling belongs to PartitionLog.}
\lessonhints{Use \method{RecordCodec.encodedSize} for each complete record, then compare with the active segment's remaining bytes.}{Walk the batch while tracking one global nextOffset and current segment size; when a read enters a new segment, carry forward the same remaining budget.}{Check directory bases, each segment's exact size, mutationVersion, and record-by-record results after reopen.}
\lessoncommand{6}
\lessonquestions{Why must an invalid batch be rejected before creating the next segment?}{What budget state must a read carry from one segment to the next?}
\paragraph*{Self-check explanation}
Prevalidation prevents an invalid request from creating an empty segment or changing the directory. A byte budget is accumulated for one call; switching physical files does not restore the bytes already used.

\lessonsection{7}{Log Retention}
\goals{Delete expired closed segments at the retained prefix while preserving the active segment, LEO, and actual new start.}{Distinguish physical retention from consumer progress.}
\background{The current retained range is \pathcode{[logStartOffset,LEO)}. deleteBefore is an explicit synchronous call; it does not run automatically when a consumer commits an offset. A segment cannot be split by individual record: only closed segments can be deleted, and the last segment must remain active so the log has an append target. The new start is the base of the first remaining segment; it may be lower than the requested threshold because segment boundaries need not match that threshold.}
\providededit{The \pathcode{PartitionLog} skeleton owns multi-segment discovery, the start/end offsets, and closing paths, and provides \method{mutationVersion()} and private \method{markMutation()}.}{Implement \method{public synchronized long deleteBefore(long offset)} in \pathcode{exercises/src/main/java/io/simplekafka/storage/PartitionLog.java}.}
\subsubsection*{Prerequisites}
Complete \stepref{6}{segmented log}; understand the offset ranges determined by a segment base and the next segment base, and why the active segment cannot be deleted. A consumer commit is separate state and is not an input to this method.
\subsubsection*{Inputs, outputs, and state}
\begin{itemize}
\item offset is a nonnegative long retention bound. A negative value reports \method{INVALID_REQUEST} and changes no state. Return the actual logStartOffset, not the requested threshold.
\item Delete only a closed segment whose endOffset (the next segment base) is at or below the threshold, together with its \pathcode{.index} at the same base. Keep the active segment and at least one segment. Do not delete unfamiliar files or directories. Deletion does not move LEO, alter remaining segment contents, or adjust consumer positions.
\item If at least one segment is actually deleted, increment mutationVersion once before the first deletion (at most once per call). If no segment qualifies, return the current start and leave the version unchanged. A read below the new start reports \method{OFFSET_OUT_OF_RANGE}; do not silently seek to the start. A closed log or underlying I/O failure reports \method{STORAGE_ERROR}; if deletion fails midway, the operation does not promise to roll back deletions that already happened.
\end{itemize}
\lessoncontract{Delete only complete closed segments as units and delete their log/index pair; never remove the active or final segment. Reopen must recover the same logStartOffset and LEO; a no-op does not increment mutationVersion.}
\subsubsection*{Hand calculation at retention boundaries}
Assume nine records at offsets 0–8, each with a null key and empty value, occupying 32 bytes; segmentBytes=96 makes segments at bases 0, 3, and 6, each with three records and 96 bytes. LEO=9, and base 6 is active. deleteBefore(4) deletes only base 0, whose end=3, so the actual start is 3. deleteBefore(9) then deletes base 3, whose end=6, so the start is 6. deleteBefore(1000) still keeps the active base-6 segment, so start remains 6. LEO stays 9 for all three calls; the third call deletes nothing and does not increment mutationVersion. The threshold 4 does not become the returned start 3 because deletion operates on whole segments.
\expected{\method{Step07Test}: \pathcode{deletesOnlyEligibleClosedSegmentsAndKeepsTheActiveSegment} checks end<=threshold, paired indexes, start/LEO, and unfamiliar files; \pathcode{retentionKeepsTheOnlyActiveSegmentWhetherEmptyOrNonempty} checks that the sole segment remains; \pathcode{deletesClosedSegmentAtItsExactEndOffset} checks deletion at the inclusive threshold boundary.}
\pitfalls{Deleting only \pathcode{.log} leaves an orphaned index; treating the threshold as “delete all offsets below it” would split a segment or delete active data. Consumer commits and log-file deletion are independent.}
\lessonhints{For each segment, write its base, exclusive end, and active status before checking end\(\le\)threshold.}{Separately consider a single empty segment and a single nonempty segment to ensure one active segment remains.}{List the .log/.index files after deletion, reopen, and read before the new start, at the boundary, and at LEO.}
\lessoncommand{7}
\lessonquestions{Why can deleteBefore(4) return logStartOffset=3 rather than 4?}{Even if a consumer has committed 9, why does that alone not trigger deletion?}
\paragraph*{Self-check explanation}
Retention uses whole segments, and base 3 is the next actual boundary. A commit records processing progress; it does not automatically delete files or prove no other consumer needs the old data.

\lessonsection{8}{Tail Truncation}
\goals{Remove a divergent log tail at nextOffset and make the log appendable again at that point.}{Understand half-open ranges, segment boundaries, and index consistency.}
\background{truncateTo(nextOffset) means “the offset to write next”: retain all records whose offset is less than nextOffset and remove the remaining tail. It is opposite to retention—retention removes an old prefix, truncation removes a newer suffix. If the target lies within a segment, truncate that file at the target record's byte boundary and rebuild its index; if it equals a segment base, keep an empty active segment at that base. This local log operation by itself does not establish a safe replica switch or change consumer positions.}
\providededit{Earlier steps established segment close, file truncation, index rebuilding, and the global segment table; \pathcode{PartitionLog} provides \method{mutationVersion()} and private \method{markMutation()}.}{Implement \method{public synchronized void truncateTo(long nextOffset)} in \pathcode{exercises/src/main/java/io/simplekafka/storage/PartitionLog.java}.}
\subsubsection*{Prerequisites}
Complete \stepref{7}{log retention}; know that actual logStartOffset may exceed zero and that LEO is an exclusive end. Truncation must stay within the retained range and cannot restore a prefix already deleted by retention.
\subsubsection*{Inputs, outputs, and state}
\begin{itemize}
\item nextOffset must be in the closed range \pathcode{[logStartOffset,LEO]}. An out-of-range value reports \method{OFFSET_OUT_OF_RANGE} and leaves files, indexes, segment table, LEO, and mutationVersion unchanged. The method returns no value.
\item On success, retain every record with offset<nextOffset; delete later records/segments and their indexes; set LEO to nextOffset. nextOffset=LEO is a no-op. A valid change increments mutationVersion once before the first file write/delete; invalid and no-op calls do not increment it. Unfamiliar files remain. A closed log or underlying I/O failure reports \method{STORAGE_ERROR}; an I/O failure after a valid change starts does not promise rollback, and the incremented mutationVersion does not revert.
\item After success, append can run immediately and its first record must receive exactly nextOffset. Reopen must recover an active segment and index consistent with the truncated log.
\end{itemize}
\lessoncontract{The truncation boundary always means the next record's offset; there is no interpretation “keep record 5 and write 5.” At a segment boundary, preserve an active segment for future appends; inside a segment, adjust the log tail and its index together.}
\subsubsection*{Hand calculation: four independent truncations}
Reuse Step 7's fixture: assume records 0–8 each occupy 32 bytes, segmentBytes=96, three full segments at bases=[0,3,6] with sizes=[96,96,96], LEO=9 and start=0. Each row below starts from that same fixture and applies one operation:
\begin{enumerate}
\item truncateTo(5): retain offsets 0–4; bases=[0,3], sizes=[96,64], LEO=5.
\item truncateTo(3): retain 0–2; bases=[0,3], sizes=[96,0]. The empty base-3 segment receives the next offset 3.
\item truncateTo(0): retain an empty log; only base 0 remains, size 0, and the next append starts at 0.
\item truncateTo(9): target equals LEO, so nothing changes; bases/sizes remain [0,3,6]/[96,96,96] and mutationVersion does not increment.
\end{enumerate}
\expected{\method{Step08Test}: \pathcode{truncatesAtBoundariesWithinSegmentsAtLeoAndAtLogStart} checks segment interiors, boundaries, LEO, and log start; \pathcode{truncateRebuildsMixedSizeSegmentStateAndIndexPoints} checks log/index consistency; \pathcode{invalidAndRepeatedTruncationsPreserveTheRetainedFiles} checks invalid and no-op calls; \pathcode{truncatesThenRetainsAcrossRestartAndContinuesAtTheReopenedEnd} checks reopen, retention, and further append; \pathcode{fixedSeedOperationsMatchAnIndependentSegmentAndOffsetModel} checks repeated state transitions.}
\pitfalls{truncateTo(5) keeps offset 4, not offset 5; 5 is the next position. Changing the log file without trimming/rebuilding its index can make a reader jump to a deleted divergent record.}
\lessonhints{Translate nextOffset into “retain offsets below it” before deciding what stays.}{Draw the active-file state for a target inside a segment, at a segment base, at LEO, and at logStart.}{Check disk files and indexes, close and reopen, then append one record and verify its offset equals the target.}
\lessoncommand{8}
\lessonquestions{After truncateTo(5), why is the next record written at 5 rather than 4?}{Why must the sparse index be handled after truncating the record-file tail?}
\paragraph*{Self-check explanation}
The half-open range retains every offset<5, so the last retained offset is 4 and the next is 5. Index entries are physical record-boundary hints; stale tail entries could send a read into deleted data rather than the current log.
